\begin{abstract}
Latent world models support visual control by predicting the consequences of candidate actions, while generative policies directly produce actions from observations and goals. Combining these capabilities raises a practical question: can a shared model support both direct control and explicit planning, and when does planning improve on its action generator?
We study this question with Fast-LeWAM, which shares a visual encoder and transformer across action generation, causal future prediction, latent path generation, and inverse dynamics.
The model can propose actions directly or generate a goal-conditioned latent path and decode it into actions. A common dynamics predictor scores both types of candidate by their predicted terminal distance to the goal.
This design enables comparisons between proposal spaces within the same trained model.
Completed experiments on an earlier action--dynamics model show that the benefits of joint training and actor-initialized planning depend strongly on the task. Under a revised evaluation protocol, action supervision improves planning on Cube across three training seeds but degrades it on Push-T.
These findings motivate separating candidate generation from candidate selection when evaluating shared world-action models.
\textbf{Working-draft status:} the four-mode model has intermediate development results; final performance, efficiency, and controlled training comparisons remain pending.
\end{abstract}
